% answers_short.tex —— 1992 年·全国卷·文科 · 简短答案速查
% 由 tools/make_short_answers.py 自 answers.tex 的 \gksolhead/\gkanswer 提取，
% 并与 1992/1992_数学真题_全国文_答案要点.md 逐题比对后生成，勿手改（改源 MD 后重跑）。

\begin{gknotice}
\textbf{使用说明}：本节只给最终答案便于快速核对，完整推导见解析版。答案为编写者独立推导，\textbf{非官方标准答案}。
\end{gknotice}

\gkpart{选择题}
\begin{center}\begin{tabular}{ccccccccccccccccccc}
\toprule
题号 & 1 & 2 & 3 & 4 & 5 & 6 & 7 & 8 & 9 & 10 & 11 & 12 & 13 & 14 & 15 & 16 & 17 & 18 \\
\midrule
答案 & A & D & D & C & D & B & B & D & D & D & B & A & C & D & D & C & A & C \\
\bottomrule
\end{tabular}\end{center}

\gkpart{填空题}
\begin{itemize}[leftmargin=2.6em]
  \item[\textbf{19.}] \(\dfrac14\)
  \item[\textbf{20.}] \(-\dfrac{\sqrt5}{5}\)
  \item[\textbf{21.}] \(x=-1\)
  \item[\textbf{22.}] \(\dfrac{15}{128}\)
  \item[\textbf{23.}] \(\dfrac{(x-2)^2}{4}-\dfrac{y^2}{12}=1\)
\end{itemize}

\gkpart{解答题}
\begin{itemize}[leftmargin=2.6em]
  \item[\textbf{24.}] \(\dfrac14\)。
  \item[\textbf{25.}] \(z=3+4\mathrm{i}\) 或 \(z=\dfrac53+4\mathrm{i}\)。
  \item[\textbf{26.}] \(\dfrac{a^3}{6}\)。
  \item[\textbf{27.}] \(A(-1,0)\)，\(C(5,-6)\)（条件结论：源内高线方程缺 \(x\)，按 \(x-2y+1=0\) 求得；见存疑登记）。
  \item[\textbf{28.}] (1) \(-\dfrac{24}{7}<d<-3\)；(2) \(S_6\) 最大
\end{itemize}
